\begin{center}\begin{tikzpicture}[font=\small,>=stealth,line join=round]
\draw[gray!40,rounded corners=3pt] (0,0) rectangle (14.2,8.5);\node at (7.1,7.9) {平移证明：把中心搬到原点，就得到奇函数};
\node at (2.9000000000000004,6.9){$f(x)=(x-2)^3+3$};\begin{scope}[shift={(0.7,2.9000000000000004)},x=1.1cm,y=.4cm]\draw[->,gray] (0,0)--(4,0);\draw[->,gray] (0,-2)--(0,8);\draw[blue,thick,domain=0.44999999999999996:3.55,samples=140] plot(\x,{(\x-2)^3+3});\fill[violet] (2,3) circle(2.4pt);\end{scope}\node at (10.100000000000001,6.9){$g(t)=f(2+t)-3=t^3$};\begin{scope}[shift={(10.100000000000001,4.1)},x=1.1cm,y=.4cm]\draw[->,gray] (-2,0)--(2,0);\draw[->,gray] (0,-5)--(0,5);\draw[blue,thick,domain=-1.55:1.55,samples=140] plot(\x,{(\x)^3});\fill[violet] (0,0) circle(2.4pt);\end{scope}\draw[->,orange!85!black,thick] (5.8,4.2)--(7.4,4.2);\node at (7.1,1.4){向左移2、向下移3；两个坐标方向分别使用相同单位长度。};
\node[text width=12.6cm,align=center] at (7.1,.65) {移动后横坐标为 $t=x-2$，高度为 $g(t)=f(x)-3$。};
\end{tikzpicture}\end{center}
